\section{Related Work}
\label{sec:related_work}

\subsection{KV-Cache Sharing and Reuse}
Based on the similarity of KV-Cache between layers, intra-model cache sharing methods~\citep{yang2024kvsharer, wu2024layer, sun2024you, brandon2024reducing, wu2025systematic} have been proposed to reuse shallow layers' KV-Cache for deeper layers to accelerate single LLM inference. 
Another research focus is to reuse a portion of KV-Cache (e.g., common prefix, reference documents) for the same model in multiple user queries~\citep{bang2023gptcache, ye2024chunkattention, yao2024cacheblend, qin2024mooncake, yang2025kvlink, ye2025kvcomm, eyuboglu2025cartridges}.
DroidSpeak~\citep{liu2024droidspeak} extends cache reuse to models fine-tuned from the same base model. 
% C2C improves performance
Unlike existing works that focus on computational efficiency through cache reuse, our approach leverages the KV-Cache as a medium for semantic transfer between LLMs.
% Generalize to different model sizes and families
Furthermore, unlike existing cache sharing methods that are restricted to only a single model or models with identical structure and size, our method supports sharing across different model families and varying model sizes. 

\subsection{Multi-LLM Systems}
% \subsection{Collaborative Multi-LLM Systems}
\xhdr{{Collaborative multi-LLM systems}}
% Strong to strong
Collaborative systems treat multiple LLMs as peers that exchange information to improve collective performance. 
Chain-of-Agents~\citep{zhang2024chain} and MetaGPT~\citep{hong2024metagpt} create sequential message flows where agents directly communicate using natural language interfaces. 
Mixture-of-Agents~\citep{wang2024moa}, DyLAN~\citep{liu2024dylan}, and Owl~\citep{hu2025owl} introduce layered communication architectures. Target LLMs aggregate messages from multiple models using voting or summarization mechanisms.
Multi-agent debate methods~\citep{estornell2024multi, liang2024encouraging, du2023improving} involve iterative communication rounds, letting LLM agents discuss and refine responses.
Recent works such as MCP~\cite{anthropic2024mcp} and A2A~\cite{google2025a2a} establish formal text protocols beyond natural language, standardizing agent interaction and tool usage in collaborative multi-LLM systems. 
% Do not require explicit token-by-token decoding to communicate
% Deeper semantic
These approaches rely on text-level interfaces, where communication requires one model to generate text token-by-token and another to ingest it as input. 
Our work explores a deeper and more efficient collaboration~\citep{gibberlink, zheng2025thought} by directly sharing internal KV-Cache representations.

% \subsection{Routing-based Multi-LLM Inference Systems}
\xhdr{Routing-based multi-LLM inference systems}
% Strong to weak
To accelerate LLM inference, several systems leverage multiple models with different capabilities and costs. 
Dynamic model selection methods~\citep{OpenAI2025GPT5, ong2024routellm, feng2024graphrouter, ning2024sot} route queries to different models with varying sizes and configurations to balance efficiency and performance. 
Token-level routing methods~\citep{zhang2024fast, shen2024learning, zheng2025citer, fu2025r2r} enable finer-grained selection, utilizing smaller models for simple token generation within the reasoning process of complex tasks. 
While these systems achieve efficiency through strategic model switching, they either completely drop context from other models, or simply rely on their own understandings of the context. 
% There is no contextual understanding transfer, and the small model generally does not have a strong an understanding as the big model. 
Without understanding sharing, smaller models cannot benefit from the richer representations already computed by larger models.